% L2: Generalised Bernoulli numbers and their congruences with the
% classical Bernoulli numbers.

This chapter introduces the generalised Bernoulli numbers
$\Bgen{n}{\chi}$ attached to a Dirichlet character $\chi$ and records
the properties of these numbers needed for the analytic formula for
$\hminus$.

\section{Definition}

We follow Washington's explicit formula
(Washington~(4.1), Diekmann~Prop.~45) rather than the
generating-function description: the explicit form is what is used
throughout, and it makes the dependence on the polynomial $B_n(X)$
manifest. Recall that $B_n(X) \in \Q[X]$ denotes the $n$-th Bernoulli
polynomial.

\begin{definition}[Generalised Bernoulli number]
  \label{def:gen-bernoulli}
  \lean{KummerCriterion.BernoulliGen}\leanok
  \uses{def:dirichlet-character}
Let $N \ge 1$, let $R$ be a commutative $\Q$-algebra and let
$\chi : (\Z/N\Z) \to R$ be a Dirichlet character (with $\chi(0) = 0$).
The \emph{generalised Bernoulli number} attached to $\chi$ in degree
$n \in \N$ is
\[
  \Bgen{n}{\chi}
  \;:=\;
  N^{\,n-1}\,\sum_{a \in \Z/N\Z} \chi(a)\,B_n\!\left(\frac{a}{N}\right)
  \;\in\; R,
\]
where $a$ is identified with its representative in
$\{0,1,\dots,N-1\}$, and $B_n(a/N) \in \Q$ is
sent to $R$ along $\Q \to R$. For the trivial character one recovers
the classical Bernoulli number: $\Bgen{n}{\mathbf{1}} = B_n$.
\end{definition}

\section{Vanishing in degree 0}

\begin{proposition}[$\Bgen{0}{\chi} = 0$ for non-trivial $\chi$]
  \label{prop:gen-bernoulli-zero}
  \lean{KummerCriterion.BernoulliGen_zero_of_ne_one}\leanok
  \uses{def:gen-bernoulli}
Let $R$ be an integral domain and let $\chi$ be a Dirichlet character
modulo $N \ge 1$ valued in $R$ with $\chi \ne \mathbf{1}$. Then
\[
  \Bgen{0}{\chi} = 0.
\]
\end{proposition}

\begin{proof}\leanok
The polynomial $B_0(X) = 1$ is the constant polynomial, so its
evaluation at $a/N$ equals $1$ and the algebra map sends it to
$1 \in R$. Substituting into the definition,
\[
  \Bgen{0}{\chi} \;=\; \sum_{a \in \Z/N\Z} \chi(a).
\]
For any non-trivial character $\chi$ of a finite abelian group, this
character sum vanishes: pick $b$ with $\chi(b) \ne 1$ and observe that
the substitution $a \mapsto a b$ multiplies the sum by $\chi(b)$, which
is only consistent with $\chi(b) \cdot \Sigma = \Sigma$ if $\Sigma = 0$,
since $1 - \chi(b)$ is a nonzero element of the integral domain $R$.
\end{proof}

\section{The first generalised Bernoulli number}

The case $n = 1$ is the one that appears in the analytic formula for
$\hminus$; we record three forms used throughout.

\begin{lemma}[Intermediate form for $\Bgen{1}{\chi}$]
  \label{lem:gen-bernoulli-one-intermediate}
  \lean{KummerCriterion.BernoulliGen_one_of_ne_one}\leanok
  \uses{def:gen-bernoulli}
Let $R$ be an integral domain and $\chi : (\Z/N\Z) \to R$ a non-trivial
Dirichlet character. Then
\[
  \Bgen{1}{\chi}
  \;=\;
  \sum_{a \in \Z/N\Z} \chi(a)\,\frac{a}{N}.
\]
\end{lemma}

\begin{proof}\leanok
Since $B_1(X) = X - \tfrac12$ and the exponent $1 - 1 = 0$ in the
definition makes the leading $N^{n-1}$ factor equal to $1$,
\[
  \Bgen{1}{\chi}
  = \sum_a \chi(a) \cdot \frac{a}{N}
  - \sum_a \chi(a) \cdot \frac{1}{2}.
\]
The second sum equals $\tfrac12 \cdot \sum_a \chi(a)$, which is zero
for non-trivial $\chi$ by the character-sum identity used in
\cref{prop:gen-bernoulli-zero}.
\end{proof}

\begin{lemma}[Cleared form of $\Bgen{1}{\chi}$]
  \label{lem:gen-bernoulli-one-cleared}
  \lean{KummerCriterion.natCast_mul_BernoulliGen_one_of_ne_one}\leanok
  \uses{lem:gen-bernoulli-one-intermediate}
For $\chi$ as in the previous lemma,
\[
  N \cdot \Bgen{1}{\chi}
  \;=\;
  \sum_{a \in \Z/N\Z} \chi(a)\,a,
\]
an identity valid in any commutative $\Q$-algebra (denominators have
been cleared).
\end{lemma}

\begin{proof}\leanok
Multiply the previous identity through by $N$ and use that the
algebra map $\Q \to R$ commutes with multiplication by $N \in \N$.
\end{proof}

\begin{proposition}[Parity vanishing at $n = 1$]
  \label{prop:gen-bernoulli-one-even}
  \lean{KummerCriterion.BernoulliGen_one_eq_zero_of_even_ne_one}\leanok
  \uses{lem:gen-bernoulli-one-cleared, def:even-odd}
Let $R$ be an integral domain in which $2 \ne 0$ and let $\chi$ be a
non-trivial \emph{even} Dirichlet character modulo $N > 1$ valued in
$R$, with $N$ invertible in $R$. Then
\[
  \Bgen{1}{\chi} = 0.
\]
\end{proposition}

\begin{proof}\leanok
Reindexing the sum in
\cref{lem:gen-bernoulli-one-cleared} by $a \mapsto -a$ and using
$\chi(-a) = \chi(-1) \chi(a) = \chi(a)$ (because $\chi$ is even),
\[
  \sum_a \chi(a)\,a \;=\; \sum_a \chi(a)\,(-a).\text{val},
\]
where $(-a).\text{val}$ denotes the canonical representative in
$\{0,\dots,N-1\}$. Adding the two expressions,
\[
  2 \sum_a \chi(a)\,a
  \;=\; \sum_a \chi(a)\,\bigl(a + (-a).\text{val}\bigr).
\]
The pairing identity in $\Z/N\Z$ states that
$a.\text{val} + (-a).\text{val}$ equals $0$ when $a = 0$ and $N$
otherwise. Substituting,
\[
  2 \sum_a \chi(a)\,a
  \;=\; N \cdot \sum_{a \ne 0} \chi(a)
  \;=\; N \cdot \Bigl(\sum_a \chi(a) - \chi(0)\Bigr)
  \;=\; 0,
\]
because $\chi(0) = 0$ and $\sum_a \chi(a) = 0$ for non-trivial $\chi$.
Since $2 \ne 0$ in $R$, the sum vanishes, and then $N \ne 0$ forces
$\Bgen{1}{\chi} = 0$ via \cref{lem:gen-bernoulli-one-cleared}.
\end{proof}

This is the special case at $n = 1$ of the general parity phenomenon.  For a
non-trivial Dirichlet character $\chi$ and $n \ge 1$ one has
\[
  \Bgen{n}{\chi} = 0
  \qquad\text{unless}\qquad
  \chi(-1) = (-1)^n.
\]

Use the defining sum and reindex by $a\mapsto -a$.  The Bernoulli
polynomials satisfy
\[
  B_n(1-X)=(-1)^n B_n(X),
\]
and the character contributes the scalar $\chi(-1)$. Hence
\[
  \Bgen{n}{\chi}=\chi(-1)(-1)^n\Bgen{n}{\chi}.
\]
If
\(\chi(-1)\ne (-1)^n\), the scalar multiplying \(\Bgen{n}{\chi}\) is not
$1$, so \(\Bgen{n}{\chi}=0\).

In particular: if $\chi$ is odd, $\Bgen{n}{\chi}$ can be non-zero only
for odd $n$; if $\chi$ is even, only for even $n$. Downstream this is
used in two places: (i) the cleared form of the analytic formula for
$\hminus$ in \cref{ch:hminus-formula} runs only over odd characters,
because their $\Bgen{1}{\chi}$ are the only non-vanishing terms; and
(ii) the odd special-value formula from \cref{ch:lvalue-negative}
depends on the same parity constraint.

\section{Congruences with classical Bernoulli numbers}

The remaining results in this chapter concern the classical Bernoulli
numbers $B_n \in \Q$ and their $p$-adic behaviour. They are needed in
two ways: to control the ordinary Bernoulli numbers below the boundary
$p-1$, and to handle the exceptional contribution of the boundary
character $\omega^{p-2}$ to the product side of the analytic formula
for $\hminus$.

\subsection*{$p$-integrality below the boundary}

\begin{theorem}[$p$-integrality of $B_k$ for $k < p - 1$]
  \label{thm:bernoulli-padicInt-below}
  \lean{KummerCriterion.bernoulli_mem_padicInt_of_lt_sub_one}\leanok
  \uses{def:gen-bernoulli}
Let $p$ be an odd prime and let $0 \le k < p - 1$. Then
$B_k \in \Q$ lies in $\Z_p$; equivalently, there exists $z \in \Z_p$
with $B_k = z$ in $\Q_p$.
\end{theorem}

\begin{proof}\leanok
Strong induction on $k$. The base case $k = 0$ is $B_0 = 1 \in \Z_p$.
For the inductive step assume $k \ge 1$ and that $B_j \in \Z_p$ for
all $j < k$. The recursion $\sum_{j=0}^{k} \binom{k+1}{j} B_j = 0$
extracts the top term, and after rearrangement,
\[
  (k + 1)\,B_k
  \;=\; -\sum_{j = 0}^{k - 1} \binom{k+1}{j}\,B_j.
\]
By the induction hypothesis each $B_j$ is a $p$-adic integer. The
binomial coefficients $\binom{k+1}{j}$ are integers, hence are also
$p$-adic integers. So the right-hand side is in $\Z_p$. Finally
$k + 1 \le p - 1 < p$ means $p \nmid (k+1)$, so $k + 1$ is a unit in
$\Z_p$ and we may divide. Hence $B_k \in \Z_p$.
\end{proof}

\begin{corollary}[Denominator coprimality]
  \label{cor:bernoulli-den-coprime}
  \lean{KummerCriterion.prime_not_dvd_bernoulli_den_of_lt_sub_one}\leanok
  \uses{thm:bernoulli-padicInt-below}
For $p$ an odd prime and $0 \le n < p - 1$,
\[
  p \nmid (B_n)_{\text{den}}.
\]
\end{corollary}

\begin{proof}\leanok
By \cref{thm:bernoulli-padicInt-below} the $p$-adic absolute value of
$B_n$ is at most $1$, so the denominator of $B_n$ (which equals
$|B_n|_p^{-1}$ raised to $v_p$ when negative) is a $p$-unit. Equivalently,
the cast $((B_n)_{\text{den}} : \Z_p)$ is a unit, which is equivalent
to $p \nmid (B_n)_{\text{den}}$.
\end{proof}

\subsection*{Von Staudt--Clausen at the boundary}

The recursion that pushes $B_k$ into $\Z_p$ for $k < p - 1$ breaks
down at $k = p - 1$: the scalar $k + 1 = p$ is no longer a $p$-unit.
The Bernoulli number $B_{p-1}$ has $p$ in its denominator, and the
classical Von Staudt--Clausen theorem pins this denominator down
exactly:

\begin{theorem}[Von Staudt--Clausen, $p$-local form at $n = p - 1$]
  \label{thm:von-staudt-clausen}
  \lean{KummerCriterion.bernoulli_pSubOne_add_inv_p_mem_padicInt}\leanok
  \uses{thm:bernoulli-padicInt-below}
For $p$ an odd prime,
\[
  B_{p-1} + \frac{1}{p} \;\in\; \Z_p.
\]
Equivalently, $p \cdot B_{p-1} \equiv -1 \pmod{p}$ as $p$-adic
integers.
\end{theorem}

\begin{proof}\leanok
Specialise the Bernoulli recursion at $n = p$:
\[
  \sum_{k = 0}^{p - 1} \binom{p}{k}\,B_k \;=\; 0.
\]
The top index $k = p - 1$ contributes $\binom{p}{p - 1} \cdot B_{p-1}
= p \cdot B_{p-1}$. Isolating this and splitting off $k = 0$
(which contributes $B_0 = 1$) gives
\[
  p \cdot B_{p-1} + 1
  \;=\; -\sum_{k = 1}^{p - 2} \binom{p}{k}\,B_k.
\]
For each \(1 \le k \le p - 2\), the binomial coefficient
\(\binom{p}{k}\) is divisible by \(p\); write
$\binom{p}{k} = p \cdot c_k$ for some $c_k \in \N$. By
\cref{thm:bernoulli-padicInt-below}, each $B_k$ for $1 \le k \le p - 2$
is a $p$-adic integer (note that $p - 2 < p - 1$). Therefore the
right-hand side is $p$ times an element of $\Z_p$, and dividing by
$p$ (cancellable in $\Q_p$) yields
\[
  B_{p - 1} + \frac{1}{p} \;=\; -\sum_{k=1}^{p-2} c_k \cdot B_k \in \Z_p,
\]
as claimed.
\end{proof}

\subsection*{The boundary character $\omega^{p-2}$}

The product side of the analytic formula for $\hminus$
(\cref{ch:hminus-formula}) runs over odd Dirichlet characters $\chi$
modulo $p$; identifying characters with powers of the Teichm\"uller
character $\omega$, every odd $\chi$ has the form $\omega^j$ for some
odd $j \in \{1, 3, \dots, p - 2\}$. The contribution
$\Bgen{1}{\omega^j}$ is $p$-integral for $j \in \{1, 3, \dots, p - 4\}$
(by the same recursion-based analysis used above, extended to
character twists); the exceptional case is $j = p - 2$ — the
``inverse Teichm\"uller'' character $\omega^{-1}$, which is the one
boundary character that sees the Von Staudt--Clausen pole.

We use the $\Q_p$-valued Teichm\"uller character
\(\omega : (\Z/p\Z)^\times \to \Q_p^\times\).

\begin{theorem}[Boundary character contribution]
  \label{thm:boundary-character-bernoulli}
  \lean{KummerCriterion.bernoulliGen_teichmuller_inverse_eq_p_sub_one_div_p_add_padicInt}
  \leanok
  \uses{lem:gen-bernoulli-one-cleared, def:teichmuller}
For $p$ an odd prime there exists $z \in \Z_p$ with
\[
  \Bgen{1}{\omega^{p-2}} \;=\; \frac{p - 1}{p} \;+\; z.
\]
In particular, $\Bgen{1}{\omega^{-1}}$ shares the $p$-adic pole of the
classical $B_{p-1}$ from \cref{thm:von-staudt-clausen}, with residue
$(p - 1)/p$.
\end{theorem}

\begin{proof}\leanok
Set $\chi := \omega^{p - 2}$ regarded as a Dirichlet character modulo
$p$ valued in $\Q_p$, with underlying $\Z_p$-valued lift
$\chi_{\Z_p} := \omega_{\Z_p}^{p - 2}$. Since $p$ is odd, $p - 2 \in
\{1, \dots, p - 2\}$ is not divisible by $p - 1$, so $\chi \ne
\mathbf{1}$; the cleared form
\cref{lem:gen-bernoulli-one-cleared} gives
\[
  p \cdot \Bgen{1}{\chi}
  \;=\; \sum_{a \in \Z/p\Z} \chi(a)\,a
  \;=:\; S \;\in\; \Z_p,
\]
the last identity because each summand is the product of a
$\Z_p$-valued character value and an integer representative.

We compute $S$ modulo $p$. For $a = 0$, the summand is $0$ because
$\chi_{\Z_p}(0) = 0$. For $a \ne 0$, the Teichm\"uller character lifts
the identity modulo $p$: $\omega_{\Z_p}(a) \equiv a \pmod{p}$. Therefore
\[
  \chi_{\Z_p}(a) \equiv a^{p - 2} \pmod{p},
\]
and hence
\[
  \chi_{\Z_p}(a) \cdot a \equiv a^{p - 1} \equiv 1 \pmod{p}
\]
by Fermat's little theorem.
Summing,
\[
  S \;\equiv\; \sum_{a \ne 0} 1 \;=\; p - 1 \pmod{p}.
\]
So $S - (p - 1) \in p \cdot \Z_p$; write $S - (p - 1) = p \cdot z$
with $z \in \Z_p$. Then
\[
  p \cdot \Bgen{1}{\chi} \;=\; S \;=\; (p - 1) + p \cdot z
  \;=\; p \cdot \Bigl(\frac{p - 1}{p} + z\Bigr),
\]
and cancelling the unit $p \in \Q_p^{\times}$ gives the claim.
\end{proof}

The remaining bridge to the analytic formula for $\hminus$ is Euler's
identity expressing $L(1 - n, \chi)$ in terms of $\Bgen{n}{\chi}$;
this is the subject of \cref{ch:lvalue-negative}.
